\section{总结与延伸}

\subsection{课程总结}

\begin{figure}[H]
\centering
\includegraphics[width=0.95\textwidth]{slides-images/slide-58.jpg}
\caption{Deep RL 总结：基础概念、Q-Learning 与 Policy Gradient 三大板块\protect\footnotemark}
\end{figure}
\footnotetext{Slides 第58页。}

本讲系统介绍了 Deep Reinforcement Learning 的核心内容，主要分为三个部分：

\textbf{一、RL 基础概念}
\begin{itemize}
    \item Agent 在 Environment 中通过 Action 进行交互
    \item State-Action pairs 构成学习的基本数据
    \item 目标是最大化折扣累积回报 $R_t = \sum \gamma^i r_i$
    \item Q-function $Q(s,a)$ 衡量状态-动作对的期望未来回报
\end{itemize}

\textbf{二、Value Learning / Q-Learning}
\begin{itemize}
    \item 用 DQN 近似 Q-function
    \item 策略通过 $\arg\max_a Q(s,a)$ 确定性推导
    \item 适用于离散动作空间
    \item 在 Atari 游戏中取得了超越人类的表现
\end{itemize}

\textbf{三、Policy Gradient}
\begin{itemize}
    \item 直接学习和优化策略函数 $\pi(s)$
    \item 适用于连续动作空间
    \item 通过 $\text{loss} = -\log P(a_t|s_t) \cdot R_t$ 训练
    \item 在自动驾驶、围棋等复杂应用中取得突破
\end{itemize}

\subsection{延伸阅读}

\begin{itemize}
    \item \textbf{DQN 原始论文}：Mnih et al., ``Human-level control through deep reinforcement learning,'' \textit{Nature}, 2015
    \item \textbf{REINFORCE 算法}：Williams, ``Simple statistical gradient-following algorithms for connectionist reinforcement learning,'' \textit{Machine Learning}, 1992
    \item \textbf{AlphaGo}：Silver et al., ``Mastering the game of Go with deep neural networks and tree search,'' \textit{Nature}, 2016
    \item \textbf{AlphaZero}：Silver et al., ``A general reinforcement learning algorithm that masters chess, shogi, and Go through self-play,'' \textit{Science}, 2018
    \item \textbf{VISTA 仿真器}：Amini et al., ``Learning Robust Control Policies for End-to-End Autonomous Driving From Data-Driven Simulation,'' \textit{IEEE RA-L}, 2020
    \item \textbf{RLHF}：Ouyang et al., ``Training language models to follow instructions with human feedback,'' \textit{NeurIPS}, 2022
    \item \textbf{PPO}：Schulman et al., ``Proximal Policy Optimization Algorithms,'' 2017------当前最流行的 Policy Gradient 变体
    \item \textbf{Actor-Critic 方法}：结合 Value Learning 和 Policy Gradient 的优势，同时学习价值函数和策略函数
    \item \textbf{Sitio del curso}: \url{https://introtodeeplearning.com}
\end{itemize}

\begin{knowledgebox}{Deep RL 的未来方向}
深度强化学习仍然是一个快速发展的领域。当前的研究热点包括：
\begin{itemize}
    \item \textbf{Sample Efficiency}：如何用更少的交互数据达到同等性能
    \item \textbf{Multi-Agent RL}：多个 Agent 同时学习和协作/竞争的场景
    \item \textbf{Offline RL}：利用已有的历史数据进行离线训练，避免在线交互的成本和风险
    \item \textbf{Foundation Models + RL}：将大规模预训练模型的知识融入 RL，加速学习
    \item \textbf{RLHF \& AI Alignment}：通过人类反馈训练更安全、更有用的 AI 系统
\end{itemize}
\end{knowledgebox}


%% --- Fin del contenido --- %%

\end{document}
